% !TeX root = ../main.tex

\section{Introduction}

The study of complex systems allows to explore the transition between the classical order and the statistical manifestations in the quantum regime. As in classical mechanics chaos is defined as the exponential sensitivity to initial conditions, the linearity in the Schrödinger equation does not allow for a direct translation of this definition to quantum mechanics. 
However, the phenomena known as 'Quantum Chaos' is manifested throughout the statistichal properties of the energy spectrum of the system.

The core of this project is the  Bohigas-Giannoni-Schmit (BGS) conjecture, which states that the energy levels of a quantum system whose classical counterpart is chaotic will exhibit the same statistical properties as the eigenvalues of a random matrix ensemble. 

To validate this hypothesis, we'll investigate the distribution of spacing between adjacent energy levels in two systems: one integrable and one chaotic. 

The first one is the two-dimensional \textbf{rectangular billiard}, which is an \textbf{integrable system}. In this system, the random matrix theory (RMT) predicts that the distribution of level spacings follows a Poisson distribution, which is given by:
\begin{equation}
P(s) = e^{-s}
\label{eq:Poisson_dist}
\end{equation}


The second system is the \textbf{stadium billiard}, which is a \textbf{chaotic system}. In this case, RMT predicts that the distribution of level spacings follows the Wigner-Dyson distribution, which is given by:
\begin{equation}
    P(s) = \frac{\pi}{2} s e^{-\frac{\pi}{4} s^2}
    \label{eq:Wigner-Dyson_dist}
\end{equation}

where $s$ is the normalized spacing between adjacent energy levels in both cases.

To integrate the time independent Schrödinger equation for these systems, we will use the \textbf{Finite Difference Method} (FDM): this method allows us to discretize the spatial domain and solve for the energy eigenvalues and eigenfunctions of the system.
